\section{Levantamiento non\'adico de la unidad de Pell}

Para \(k\ge1\), sea el anillo cuadr\'atico no ramificado
\begin{equation}
\mathcal R_k=(\Z/3^k\Z)[r]/(r^2-2),
\label{eq:pell-ring}
\end{equation}
y sea

\[
u=3+2r.
\]

La conjugaci\'on \(r\mapsto-r\) da
\[
N_{\mathcal R_k/(\Z/3^k\Z)}(u)
=(3+2r)(3-2r)=9-8=1,
\]
de modo que \(u\) es una unidad de norma uno. Adem\'as,
\[
u^2=17+12r,\qquad
u^4-1=576+408r=3(192+136r).
\]
El segundo factor es unidad m\'odulo tres porque reduce a \(r\ne0\). Por
tanto la valuaci\'on \(3\)-\'adica satisface
\(v_3(u^4-1)=1\). El lema de elevaci\'on para una unidad congruente con uno
m\'odulo tres produce, para todo \(j\ge0\),
\[
v_3\!\left(u^{4\cdot3^j}-1\right)=j+1.
\]
Como \(u\bmod3=2r\) y \((2r)^2=-1\), su orden m\'odulo tres es cuatro.
En consecuencia,

\begin{equation}
\operatorname{ord}(u)=4\cdot3^{k-1},
\label{eq:pell-order}
\end{equation}

y el cierre pro-finito de los grupos c\'iclicos compatibles es
\begin{equation}
\varprojlim_k\langle u\bmod3^k\rangle
\simeq C_4\times\mathbb Z_3.
\label{eq:pell-profinite}
\end{equation}
Su realizaci\'on arquimediana es

\[
3+2\sqrt2=(1+\sqrt2)^2.
\]

Esta construcci\'on demuestra la conservaci\'on de la profundidad non\'adica
por una unidad hiperb\'olica. Las unidades \(3+2\sqrt2\),
\(2+\sqrt3\) y \(30+\sqrt{899}\) forman una familia de Pell con
procedencias distintas; cualquier identificaci\'on entre ellas requiere
mapas expl\'icitos entre las extensiones correspondientes.

